\section{Krylov Reduced Deflation for MultiSignal}

Let
\[
A = I + \lambda_L L + \lambda_\theta \Theta .
\]
The first E-step solves
\[
A x_1 = y .
\]
For \(k \ge 2\), the deflated modes in Algorithm~\ref{alg:multi_mode_glasso}
solve
\[
P_{k-1} A P_{k-1} z_k = P_{k-1} y,
\qquad
x_k = P_{k-1}z_k,
\]
where
\[
P_{k-1}=I-Q_{k-1}Q_{k-1}^{\top},
\qquad
Q_{k-1}=[q_1,\ldots,q_{k-1}],
\qquad
q_i=\frac{x_i}{\|x_i\|_2}.
\]

\subsection{Krylov Reduction}

Construct the Krylov subspace
\[
\mathcal K_m(A,y)
=
\operatorname{span}\{y,Ay,A^2y,\ldots,A^{m-1}y\}.
\]
Using Lanczos for symmetric \(A\), obtain an orthonormal basis \(V_m\)
and a small tridiagonal matrix \(T_m\) such that
\[
V_m^{\top}V_m=I,
\qquad
V_m^{\top}AV_m=T_m,
\qquad
V_m^{\top}y=\beta e_1,
\]
with \(\beta=\|y\|_2\).
Approximate each mode by
\[
x_k \approx V_m s_k .
\]
If \(q_i=V_m h_i\), the original projection becomes the small-space
projection
\[
P^H_{k-1}=I-H_{k-1}H_{k-1}^{\top},
\qquad
H_{k-1}=[h_1,\ldots,h_{k-1}].
\]
The projected large system is reduced to
\[
P^H_{k-1}T_mP^H_{k-1}r_k
=
P^H_{k-1}\beta e_1,
\qquad
s_k=P^H_{k-1}r_k,
\qquad
x_k=V_m s_k .
\]

\begin{algorithm}[H]
\caption{\(m\)-step Lanczos construction using an operator}
\label{alg:operator_lanczos}
\begin{algorithmic}[1]
\Require Linear operator \(A(\cdot)\), start vector \(y\), Krylov dimension \(m\),
tolerance \(\varepsilon\)
\Ensure Krylov basis \(V_m=[v_1,\ldots,v_m]\), tridiagonal matrix \(T_m\),
and \(\beta=\|y\|_2\)
\State \(\beta\gets\|y\|_2\)
\If{\(\beta\le\varepsilon\)}
    \State \Return empty basis, empty \(T_m\), and \(\beta\)
\EndIf
\State \(v_1\gets y/\beta\), \quad \(v_0\gets 0\), \quad \(\gamma_0\gets 0\)
\For{\(j=1,\ldots,m\)}
    \State \(w\gets A(v_j)-\gamma_{j-1}v_{j-1}\)
    \State \(\alpha_j\gets v_j^{\top}w\)
    \State \(w\gets w-\alpha_jv_j\)
    \State Reorthogonalize \(w\gets w-V_j(V_j^{\top}w)\), where
    \(V_j=[v_1,\ldots,v_j]\)
    \State \(\gamma_j\gets\|w\|_2\)
    \If{\(\gamma_j\le\varepsilon\) or \(j=m\)}
        \State Set final dimension \(m'\gets j\)
        \State \textbf{break}
    \EndIf
    \State \(v_{j+1}\gets w/\gamma_j\)
\EndFor
\State Form \(V_{m'}=[v_1,\ldots,v_{m'}]\)
\State Form tridiagonal \(T_{m'}\) with diagonal
\(\alpha_1,\ldots,\alpha_{m'}\) and off-diagonal
\(\gamma_1,\ldots,\gamma_{m'-1}\)
\State \Return \(V_{m'},T_{m'},\beta\)
\end{algorithmic}
\end{algorithm}

\begin{algorithm}[H]
\caption{Projected Krylov subspace solve for one deflated mode}
\label{alg:projected_krylov_solve}
\begin{algorithmic}[1]
\Require Reduced operator \(T_m\), reduced right-hand side \(b_m=\beta e_1\),
current reduced orthonormal basis \(H_{k-1}\)
\Ensure Reduced mode coordinate \(s_k\)
\State Define \(P^H_{k-1}(v)\gets v-H_{k-1}(H_{k-1}^{\top}v)\)
\State Define the projected reduced operator
\[
\mathcal T_{k-1}(v)\gets P^H_{k-1}T_mP^H_{k-1}v
\]
\State \(b_k\gets P^H_{k-1}b_m\)
\State Solve \(\mathcal T_{k-1}(r_k)=b_k\) in \(\mathbb R^m\)
\State \(s_k\gets P^H_{k-1}r_k\)
\State \Return \(s_k\)
\end{algorithmic}
\end{algorithm}

\begin{algorithm}[H]
\caption{Krylov reduced multi-mode signal generation}
\label{alg:krylov_reduced_multimode}
\begin{algorithmic}[1]
\Require Observation \(y\), operators \(L,\Theta\), weights
\(\lambda_L,\lambda_\theta\), modes \(K\), Krylov dimension \(m\)
\Ensure Mode matrix \(X=[x_1,\ldots,x_K]\)
\State Define \(A(v)\gets v+\lambda_L L(v)+\lambda_\theta\Theta(v)\)
\State Run Algorithm~\ref{alg:operator_lanczos} on \(A\) with start vector \(y\)
\State Obtain \(V_m,T_m,\beta\), where \(V_m^{\top}y=\beta e_1\)
\State Solve \(T_m s_1=\beta e_1\)
\State \(x_1\gets V_m s_1\)
\State \(h_1\gets s_1/\|s_1\|_2\), \quad \(H_1\gets[h_1]\), \quad \(X\gets[x_1]\)
\For{\(k=2,\ldots,K\)}
    \State Run Algorithm~\ref{alg:projected_krylov_solve} with
    \(T_m\), \(b_m=\beta e_1\), and \(H_{k-1}\) to obtain \(s_k\)
    \State \(x_k\gets V_m s_k\)
    \State \(h_k\gets s_k/\|s_k\|_2\)
    \State \(H_k\gets[H_{k-1},h_k]\), \quad \(X\gets[X,x_k]\)
\EndFor
\State \Return \(X\)
\end{algorithmic}
\end{algorithm}

\subsection{Complexity}

The direct projected CG implementation costs approximately
\[
O\left(\sum_{k=2}^{K}t_k\cdot \operatorname{cost}(A\text{-apply})\right),
\]
where \(t_k\) is the number of CG iterations for the \(k\)-th deflated mode.
The Krylov reduced method costs approximately
\[
O\left(m\cdot\operatorname{cost}(A\text{-apply})\right)
+O(Km^3)+O(Kmn).
\]
Here \(O(Km^3)\) comes from solving the small projected systems, and
\(O(Kmn)\) comes from reconstructing \(x_k=V_m s_k\).
When \(m\ll n\), the expensive applications of \(A\) are paid mainly during
the Lanczos phase rather than repeatedly inside each deflated CG solve.

\subsection{Remarks}

The reduced method is exact only when the Krylov subspace captures the
required solution space; otherwise it is an approximation controlled by
\(m\). Larger \(m\) improves accuracy at the cost of memory and small-system
work. If \(A\) is ill-conditioned, a preconditioned Krylov method may be
needed.
